\begin{claim}{C.45}{Proved}{Two actual boundary quotients give explicit planar coordinates}
Let $\Phi_1:O\to B_1$ and $\Phi_2:O\to B_2$
be two actual nontrivial boundary quotients
of the same open orbit as in C.42.
Use the same source universal covering
$p:X\to O$ and the real target coverings
$q_i:\R\to B_i$, and choose lifts
$q_iF_i=\Phi_i p$.
Exactly one of the following alternatives holds.
Either there is a homeomorphism
$f:\R\to\R$ with
\[
                         F_2=fF_1,                     \tag{C.45.1}
\]
or the pair $(F_1,F_2)$ is a homeomorphism
onto an open domain $\Omega\subset\R^2$
whose horizontal and vertical sections
are intervals and whose two projections
are onto $\R$.
In the second alternative, explicit
independent changes of the two real
coordinates, and a possible exchange
of them, carry $\Omega$ onto one of
\[
 \R^2,\qquad
 \{(x,y):x<y\},\qquad
 \{(x,y):x<y<x+1\}.                                    \tag{C.45.2}
\]
The full increasing coordinate-pair
groups preserving these three domains
are, respectively,
\[
 \operatorname{Homeo}^+(\R)^2,\qquad
 \{(g,g):g\in\operatorname{Homeo}^+(\R)\},\qquad
 \{(g,g):g\in\operatorname{Homeo}^+(\R),\
                         g(x+1)=g(x)+1\}.               \tag{C.45.3}
\]
The actual compatible lifted action
maps into the corresponding group.
These statements retain both original
covering equations and make no
identification of their downstairs
component labels.
\end{claim}

\begin{proof}
\emph{A common compatible group.}
Use the closed subgroup
\[
\begin{split}
 \mathcal P=\{(l,a,b_1,b_2)\in
 L\times\operatorname{Homeo}(X)
       \times\operatorname{Homeo}^+(\R)^2:\qquad&\\
 pa&=(l|_O)p,\\
 q_i b_i&=(l|_{B_i})q_i,\qquad
 F_i a=b_iF_i\quad(i=1,2)\}.
\end{split}                                             \tag{C.45.4}
\]
The equations are closed in the
product of the indicated Polish
groups. For every lift $a$ of
$l|_O$, each target lift $b_i$
exists uniquely after its deck
translation is chosen at one point,
exactly as in C.42. Consequently
$\mathcal P$ acts transitively on
the source covering plane.
Both $F_i$ are open onto $\R$
with connected proper-line fibers,
by C.43.

Fix $t\in\R$ and set
$C=F_1^{-1}(t)$.
The closed subgroup of $\mathcal P$
with $b_1(t)=t$ acts transitively
on this whole line. Its action
on the interval $I=F_2(C)$ is
also transitive, by the equality
$F_2a=b_2F_2$. All these target
maps are increasing.
A nondegenerate interval admitting
such a transitive action has no
included endpoint, so it is an
open interval.
The Polish open orbit-map theorem,
applied to the actions on $C$
and $I$, makes
$F_2|_C:C\to I$ the open map
between their homogeneous-space
quotients: the point stabilizer
in $C$ is contained in the
stabilizer of its image in $I$.
It is therefore a strictly
monotone homeomorphism.

If instead $I$ is a point,
transitivity on the source
plane transports this constancy
to every $F_1$-fiber. There is
then a function $f$ satisfying
(C.45.1). Since $F_1$ is an
open quotient, $f$ is continuous.
For every open $W\subset\R$,
\[
              f(W)=F_2(F_1^{-1}(W))
\]
is open. Thus $f$ is open and
onto, hence a homeomorphism
of the line. This is the first
alternative.

In the nonconstant case,
$(F_1,F_2)$ is injective,
since $F_2$ is injective on
each $F_1$-fiber. Invariance
of domain makes it a
homeomorphism onto an open
domain $\Omega$.
Each vertical section is the
open interval $I$ just proved.
Interchanging the two indices
gives the horizontal statement;
the constant alternative in
that direction would contradict
the already proved two-dimensional
open image.
Surjectivity of the two $F_i$
gives both full coordinate
projections.

\emph{The endpoint maps are actual homogeneous quotients.}
The coordinate action
$(b_1,b_2)$ of $\mathcal P$
preserves $\Omega$ and is
transitive on it.
Write its vertical intervals as
\[
 \Omega_x=(\alpha(x),\beta(x)),
 \quad -\infty\leq\alpha(x)<\beta(x)\leq+\infty.
                                                               \tag{C.45.5}
\]
The two endpoints satisfy the
exact equations
\[
 \alpha(b_1x)=b_2(\alpha(x)),\qquad
 \beta(b_1x)=b_2(\beta(x)),                               \tag{C.45.6}
\]
where increasing homeomorphisms
fix the appropriate infinite
endpoints.
Transitivity of the first
coordinate action implies that
$\alpha$ is either everywhere
finite or identically $-\infty$,
and that $\beta$ is either
everywhere finite or identically
$+\infty$.

Each finite endpoint map is a
homeomorphism of the real line.
For example, choose $x_0$ with
$\alpha(x_0)\in\R$. The formula
\[
 \alpha(b_1x_0)=b_2(\alpha(x_0))
\]
descends the continuous second
evaluation through the open
first orbit map
$\mathcal P\to\R$.
It proves continuity of
$\alpha$ in its actual
topology. Its image is
nonempty and invariant under
the transitive second coordinate
action, hence is all $\R$.
The stabilizer of $x_0$ in
the first action fixes
$\alpha(x_0)$ in the second
by (C.45.6). The two open
orbit maps identify $\alpha$
with the associated open
coset-space map.
It follows that $\alpha$ is
continuous, open, and onto,
and therefore strictly monotone
and a homeomorphism.
The same proof applies to
finite $\beta$.
This verifies the endpoint
topology without assuming
regularity of a boundary graph.

\emph{The three domains and the coordinate changes.}
If both endpoints are infinite,
$\Omega=\R^2$.
If precisely one endpoint is
finite, use its inverse as
the second coordinate.
The defining inequality becomes
$x<y$ or $y<x$, according to
which endpoint and monotonicity
occur. In the latter case
exchange the two coordinates.

Suppose both endpoints are
finite. They have the same
monotonicity: an increasing
homeomorphism and a decreasing
homeomorphism cannot satisfy
$\alpha(x)<\beta(x)$ for all
real $x$, by their limits at
the two ends.
When both increase, replace
$y$ by $\alpha^{-1}(y)$ and put
$\tau=\alpha^{-1}\beta$.
When both decrease, replace
$y$ by $\beta^{-1}(y)$ and put
$\tau=\beta^{-1}\alpha$.
In either case the domain is
now exactly
\[
                        x<y<\tau(x),
 \qquad \tau\in\operatorname{Homeo}^+(\R),\
                          \tau(x)>x.                    \tag{C.45.7}
\]
Choose $x_0$ and any increasing
homeomorphism
$h:[x_0,\tau(x_0)]\to[0,1]$.
Extend it by
$h(\tau^n x)=h(x)+n$.
The iterates of $\tau$ exhaust
the two ends, since $\tau$
has no fixed point, and the
endpoint formulas agree.
Thus $h$ is an increasing
homeomorphism of $\R$ with
$h\tau h^{-1}(x)=x+1$.
Applying $h$ to both
coordinates in (C.45.7)
gives the final domain in
(C.45.2).
These are the asserted
explicit source coordinate
maps; a reversal in one
coordinate still conjugates
increasing coordinate actions
to increasing actions.

Finally, an increasing
coordinate pair preserving
$x<y$ must preserve its
boundary graph $y=x$,
which says $g=h$.
For $x<y<x+1$, its lower
and upper endpoint graphs
are separately preserved,
which says
$h(x)=g(x)$ and
$h(x+1)=g(x)+1$.
Conversely these identities
preserve the indicated
domains exactly.
The unrestricted plane case
has no additional equation.
This proves (C.45.3).

All induced actions used
above arise from (C.45.4).
An element acting trivially
on the covering plane may
still act elsewhere under
its $l$-coordinate; no
faithfulness of that global
coordinate has been imposed.
The original equations
$q_iF_i=\Phi_i p$ remain
in force, so the separate
deck shifts and component
labels of C.44 are retained.
\end{proof}

\paragraph{Scope of the reduction.}
C.45 gives the relation between
two specified actual boundary
maps, including all coordinate
changes. A third quotient can
impose an additional equation
on the group in (C.45.3).
Preserving all such equations
and obtaining a global
frontier-compatible contraction
are not consequences of this
pairwise classification.
